\chapter{2021年上海卷（春）}

\section{填空题}

\begin{problem}
等差数列 \(\{a_n\}\) 中，\(a_1=3\)，\(d=2\)，则 \(a_{10}=\fillinblank{}\).
\end{problem}

\begin{answer}
\(21\)
\end{answer}

\begin{solution}
由等差数列通项公式\(a_n=a_1+\left( n-1 \right)d\)，将\(a_1=3\)，\(d=2\)，\(n=10\)代入得\(a_{10}=3+9\times 2=21\)，回代满足公差为\(2\)的等差关系，故填\(21\).
\end{solution}

\begin{problem}
已知复数 \(z=1-3\i\)，则 \(|\bar z-\i|=\fillinblank{}\).
\end{problem}

\begin{answer}
\(\sqrt5\)
\end{answer}

\begin{solution}
复数\(z=1-3\i\)的共轭复数为\(\bar z=1+3\i\)，于是\(\bar z-\i=1+2\i\)，其模为\(\left| \bar z-\i \right|=\sqrt{1^2+2^2}=\sqrt{5}\)，平方回代\(1^2+2^2=5\)相符，故填\(\sqrt{5}\).
\end{solution}

\begin{problem}
不等式 \(\dfrac{2x+5}{x-2}<1\) 的解集为\fillinblank{}.
\end{problem}

\begin{answer}
\((-7,2)\)
\end{answer}

\begin{solution}
不等式定义域要求\(x\ne 2\)，移项通分得\(\dfrac{2x+5}{x-2}-1=\dfrac{x+7}{x-2}<0\)，该分式小于零当且仅当分子分母异号，即\(x+7\)与\(x-2\)异号，解得\(-7<x<2\)，端点处分式为零或无意义均舍去，故解集为\(\left( -7,2 \right)\).
\end{solution}

\begin{problem}
已知圆柱的底面半径为 \(1\)，母线长为 \(2\)，则其侧面积为\fillinblank{}.
\end{problem}

\begin{answer}
\(4\pi\)
\end{answer}

\begin{solution}
圆柱侧面展开为长\(2\pi r\)、宽\(l\)的矩形，故侧面积公式为\(S=2\pi rl\)，将\(r=1\)，\(l=2\)代入得\(S=2\pi\times 1\times 2=4\pi\)，量纲为面积单位，故填\(4\pi\).
\end{solution}

\begin{problem}
直线 \(x=-2\) 与直线 \(\sqrt3x-y+1=0\) 的夹角为\fillinblank{}.
\end{problem}

\begin{answer}
\(\dfrac{\pi}{6}\)
\end{answer}

\begin{solution}
直线\(x=-2\)垂直于\(x\)轴，倾斜角为\(\dfrac{\pi}{2}\)，直线\(\sqrt{3}x-y+1=0\)化为\(y=\sqrt{3}x+1\)，斜率为\(\sqrt{3}\)，倾斜角为\(\dfrac{\pi}{3}\)，两倾斜角之差的绝对值为\(\dfrac{\pi}{6}\)，落在直线夹角范围\(\left[ 0,\dfrac{\pi}{2} \right]\)内，故夹角为\(\dfrac{\pi}{6}\).
\end{solution}

\begin{problem}
方程组 \(\begin{cases}
a_1x+b_1y=c_1,\\
a_2x+b_2y=c_2
\end{cases}\) 无解，则 \(\begin{vmatrix}a_1&b_1\\a_2&b_2\end{vmatrix}=\fillinblank{}\).
\end{problem}

\begin{answer}
\(0\)
\end{answer}

\begin{solution}
记系数行列式\(D=\begin{vmatrix}a_1 & b_1\\a_2 & b_2\end{vmatrix}\)，若\(D\ne 0\)，由克拉默法则方程组有唯一解，与已知无解矛盾，所以必有\(D=0\)，该结论对两直线平行、重合及系数退化情形均成立，故填\(0\).
\end{solution}

\begin{problem}
\((x+1)^n\) 的二项展开式中有且仅有 \(x^3\) 的系数为最大值，则 \(x^3\) 的系数为\fillinblank{}.
\end{problem}

\begin{answer}
\(20\)
\end{answer}

\begin{solution}
\(\left( x+1 \right)^n\)的通项为\(T_{k+1}=\mathrm C_n^k x^k\)，其中\(k=0,1,\cdots,n\)为整数，故\(x^k\)的系数为\(\mathrm C_n^k\)，若\(n<3\)则\(x^3\)系数为\(0\)而常数项为\(1\)，不可能是唯一最大值，故\(n\ge 3\)，由比值\(\mathrm C_n^{k+1}/\mathrm C_n^k=\left( n-k \right)/\left( k+1 \right)>1\)当且仅当\(k<\left( n-1 \right)/2\)知二项式系数先严格增后严格减，若\(n=2m+1\)为奇数则\(\mathrm C_n^m=\mathrm C_n^{m+1}\)并列最大，唯一性不成立，若\(n=2m\)为偶数则唯一最大值在\(k=m\)处，由\(m=3\)得\(n=6\)，此时系数列为\(1\)，\(6\)，\(15\)，\(20\)，\(15\)，\(6\)，\(1\)，确以\(x^3\)系数独占最大，故\(x^3\)的系数为\(\mathrm C_6^3=20\).
\end{solution}

\begin{problem}
已知函数 \(f(x)=3^x+\dfrac{a}{3^x+1}\)（\(a>0\)） 的最小值为 \(5\)，则 \(a=\fillinblank{}\).
\end{problem}

\begin{answer}
\(9\)
\end{answer}

\begin{solution}
指数函数\(t=3^x\)在\(\R\)上连续严格递增，值域为\(t>0\)，原函数最值等价于\(g\left( t \right)=t+\dfrac{a}{t+1}\)在\(t>0\)上的最小值，\(g'\left( t \right)=1-\dfrac{a}{\left( t+1 \right)^2}\)，若\(0<a\le 1\)则\(\left( t+1 \right)^2>1\ge a\)，\(g'\left( t \right)>0\)恒成立，\(g\)在\(\left( 0,+\infty \right)\)上严格递增，对任一\(t\)取\(t/2\)更小即无最小值，与题设矛盾，故必\(a>1\)，此时唯一驻点为\(t+1=\sqrt{a}\)，即\(t_0=\sqrt{a}-1>0\)，左减右增为全局最小点，最小值为\(g\left( t_0 \right)=\left( \sqrt{a}-1 \right)+\sqrt{a}=2\sqrt{a}-1\)，令其等于\(5\)得\(\sqrt{a}=3\)，即\(a=9\)，回代\(t_0=2\)，即\(x_0=\log_3 2\)处\(f=2+9/3=5\)，且\(\left( t+1 \right)+9/\left( t+1 \right)-1\ge 2\times 3-1=5\)取等可达，故最小值确为\(5\)，所以\(a=9\).
\end{solution}

\begin{problem}
在无穷等比数列 \(\{a_n\}\) 中，\(\lim\limits_{n\to\infty}(a_1-a_n)=4\)，则 \(a_2\) 的取值范围为\fillinblank{}.
\end{problem}

\begin{answer}
\((-4,0)\cup(0,4)\)
\end{answer}

\begin{solution}
设\(a_n=a_1q^{n-1}\)，其中\(a_1\ne 0\)，\(q\ne 0\)，否则从第二项起相邻两项的比值无意义，不合等比数列的定义，于是排除\(q=0\)的情形. 若\(q=1\)，则\(a_1-a_n=0\)，极限为\(0\)，与题设\(4\)矛盾. 若\(q=-1\)，则\(a_1-a_n\)在\(0\)与\(2a_1\)之间交替，不收敛. 若\(\left| q \right|>1\)，则\(\left| a_n \right|\)趋于无穷大，\(a_1-a_n\)发散. 于是极限存在并等于\(4\)要求\(\left| q \right|<1\)，此时\(\displaystyle\lim_{n\to\infty}a_n=0\)，由题设\(a_1-0=4\)，得到\(a_1=4\). 反之，当\(a_1=4\)且\(-1<q<1\)，\(q\ne 0\)时，数列为无穷等比数列且极限条件成立. 此时\(a_2=4q\)，\(q\)取遍\(\left( -1,0 \right)\)，\(\left( 0,1 \right)\)，故\(a_2\)的取值范围为\(\left( -4,0 \right)\cup\left( 0,4 \right)\).
\end{solution}

\begin{problem}
某人某天需要运动总时长大于等于 \(60\) 分钟，现有五项运动可选，时间分别为 \(30\)，\(20\)，\(40\)，\(30\)，\(30\) 分钟，则有\fillinblank{} 种运动组合.
\end{problem}

\begin{answer}
\(23\)
\end{answer}

\begin{solution}
按本页现稿题干（五项时长分别为\(30\)，\(20\)，\(40\)，\(30\)，\(30\)分钟，总时长至少\(60\)分钟）与原PDF第2页表格（含7点—12点时段行及时长行）计数条件一致，结论相同. 五项运动的每项只有选或不选两种状态，去掉空选后共有\(2^5-1=31\)种非空组合. 总时长小于\(60\)分钟的组合逐类清点如下. 单项组合有\(A\)，\(B\)，\(C\)，\(D\)，\(E\)共\(5\)种，时长分别为\(30\)，\(20\)，\(40\)，\(30\)，\(30\)分钟，均小于\(60\)分钟. 两项组合中，含\(20\)分钟项与一项\(30\)分钟项的和为\(50\)分钟，共有\(B+A\)，\(B+D\)，\(B+E\)共\(3\)种，而\(20+40=60\)，\(30+30=60\)，其余两项和均超过\(60\)分钟，因此小于\(60\)分钟的两项组合恰为上述\(3\)种. 任何三项组合的时长至少为\(20+30+30=80\)分钟，四项与五项组合时长更大，因此三项及以上的组合时长都不小于\(60\)分钟. 于是总时长小于\(60\)分钟的组合共\(5+3=8\)种，总时长大于等于\(60\)分钟的组合有\(31-8=23\)种.
\end{solution}

\begin{problem}
已知椭圆 \(x^2+\dfrac{y^2}{b^2}=1\)（\(0<b<1\)） 的左，右焦点为 \(F_1\)，\(F_2\)，以 \(O\) 为顶点，\(F_2\) 为焦点作抛物线交椭圆于 \(P\)，且 \(\angle PF_1F_2=45^\circ\)，则抛物线的准线方程是\fillinblank{}.
\end{problem}

\begin{answer}
\(x=1-\sqrt2\)
\end{answer}

\begin{solution}
椭圆\(x^2+\dfrac{y^2}{b^2}=1\)中\(0<b<1\)，故长轴在\(x\)轴上，记\(c=\sqrt{1-b^2}\in\left( 0,1 \right)\)，则\(F_1=\left( -c,0 \right)\)，\(F_2=\left( c,0 \right)\). 以原点为顶点、\(F_2\)为焦点的抛物线开口向右，方程为\(y^2=4cx\)，准线为\(x=-c\). 设交点\(P=\left( x_0,y_0 \right)\)满足\(y_0^2=4cx_0\)与\(x_0^2+\dfrac{y_0^2}{b^2}=1\)，且\(\angle PF_1F_2=45^\circ\). 直线\(F_1F_2\)沿\(x\)轴正方向，而\(\overrightarrow{F_1P}=\left( x_0+c,y_0 \right)\)，夹角为\(45^\circ\)推出\(\dfrac{\left| y_0 \right|}{x_0+c}=\tan 45^\circ=1\)，即\(\left| y_0 \right|=x_0+c\). 由对称性取\(y_0>0\)的情形，另一侧符号不影响\(x_0\)与\(c\)的方程，此时\(y_0=x_0+c\). 代入抛物线方程得到\(\left( x_0+c \right)^2=4cx_0\)，整理为\(\left( x_0-c \right)^2=0\)，于是\(x_0=c\)，\(y_0=2c\). 再代入椭圆方程，并用\(b^2=1-c^2\)消去\(b\)，得到\(c^2+\dfrac{4c^2}{1-c^2}=1\). 记\(u=c^2\)，上式化为\(u+\dfrac{4u}{1-u}=1\)，通分整理为\(u^2-6u+1=0\)，解为\(u=3\pm 2\sqrt{2}\). 因为\(0<u<1\)，舍去大于\(1\)的根\(3+2\sqrt{2}\)，保留\(u=3-2\sqrt{2}=\left( \sqrt{2}-1 \right)^2\)，开方并取正值得到\(c=\sqrt{2}-1\). 回代检验\(x_0+c=2c>0\)，\(P\)同时满足椭圆与抛物线方程，夹角确为\(45^\circ\)，相符. 故抛物线的准线方程为\(x=-c=1-\sqrt{2}\).
\end{solution}

\begin{problem}
已知 \(\theta>0\)，对任意 \(n\in\mathbb N^*\)，总存在实数 \(\varphi\)，使得 \(\cos(n\theta+\varphi)<\dfrac{\sqrt3}{2}\)，则 \(\theta\) 的最小值是\fillinblank{}.
\end{problem}

\begin{answer}
\(\dfrac{2\pi}{5}\)
\end{answer}

\begin{solution}
记\(I=\left[ -\dfrac{\pi}{6},\dfrac{\pi}{6} \right]\)为单位圆上余弦不小于\(\dfrac{\sqrt{3}}{2}\)的闭弧，其长度为\(\dfrac{\pi}{3}\)，其余开弧\(J=\left( \dfrac{\pi}{6},2\pi-\dfrac{\pi}{6} \right)\)的长度为\(\dfrac{5\pi}{3}\). 题意指存在实数\(\varphi\)，对一切正整数\(n\)有\(\cos\left( n\theta+\varphi \right)<\dfrac{\sqrt{3}}{2}\)，即数列\(\beta_n=n\theta+\varphi\)按模\(2\pi\)恒落在\(J\)中. 当\(\theta=\dfrac{2\pi}{5}\)时，取\(\varphi=\pi\)，则\(\beta_n\)按模\(2\pi\)只取\(\pi\)，\(\dfrac{7\pi}{5}\)，\(\dfrac{9\pi}{5}\)，\(\dfrac{\pi}{5}\)，\(\dfrac{3\pi}{5}\)五个值. 其中\(\dfrac{\pi}{5}=36^\circ>30^\circ\)，而余弦在\(\left[ 0,\pi \right]\)上严格递减，所以\(\cos\dfrac{\pi}{5}<\cos\dfrac{\pi}{6}=\dfrac{\sqrt{3}}{2}\)，其余四点的余弦经对称化简后同样小于\(\dfrac{\sqrt{3}}{2}\)，因此\(\theta=\dfrac{2\pi}{5}\)满足要求. 下证任何\(0<\theta<\dfrac{2\pi}{5}\)都不满足要求. 假设相反，即存在\(\varphi\)使所有\(\beta_n\)按模\(2\pi\)落在\(J\)中. 若\(0<\theta\le\dfrac{\pi}{3}\)，把\(\beta_n\)提升为实数列\(x_n=n\theta+\varphi\)，它每次向前移动\(\theta\)，而禁弧\(I\)的长度为\(\dfrac{\pi}{3}\ge\theta\). 若某一步从容许区间跨出，则起点在禁弧上边界之上不到\(\theta\)处，落点必定掉进长度不小于步长的禁弧内部，与恒落在\(J\)中矛盾，而\(x_n\)趋于正无穷不可能永远停留在单个长度为\(\dfrac{5\pi}{3}\)的区间提升副本中，矛盾. 若\(\dfrac{\pi}{3}<\theta<\dfrac{2\pi}{5}\)，考察子列\(\gamma_m=\beta_{n_0+5m}\)，它每次按模\(2\pi\)向前移动\(5\theta\)，等价于向后移动\(\delta=2\pi-5\theta\in\left( 0,\dfrac{\pi}{3} \right)\). 把该子列提升为每次严格减少\(\delta\)的实数列，它趋于负无穷，不可能永远停留在容许区间的提升副本中，而每次向后跨越禁弧时，因步长\(\delta\)小于禁弧长度\(\dfrac{\pi}{3}\)，跨越必定落入禁弧内部，同样与恒落在\(J\)中矛盾. 两种情形都导出矛盾，所以\(0<\theta<\dfrac{2\pi}{5}\)时不存在满足要求的\(\varphi\). 结合\(\theta=\dfrac{2\pi}{5}\)的例子，\(\theta\)的最小值为\(\dfrac{2\pi}{5}\).
\end{solution}
\section{单选题}

\begin{problem}
下列函数中，在定义域内存在反函数的是（\quad）.
\choices{\(x^2\)}{\(\sin x\)}{\(2^x\)}{\(y=1\)}
\end{problem}

\begin{answer}
C
\end{answer}

\begin{solution}
在定义域内存在反函数等价于函数在定义域上是一一对应. \(y=x^2\)在\(\R\)上不是单调的，\(1\)与\(-1\)取相同的函数值，故不是一一对应，排除A. \(y=\sin x\)在\(\R\)上是周期函数，不同自变量取相同的函数值，故不是一一对应，排除B. 常值函数\(y=1\)把一切自变量映为\(1\)，不是一一对应，排除D. \(y=2^x\)在\(\R\)上严格递增，不同自变量对应不同的函数值，且值域上的每一点都有原像，因此是一一对应，存在反函数，故选C.
\end{solution}

\begin{problem}
已知集合 \(A=\{x\mid x^2-x-2\ge0\}\)，\(B=\{x\mid x>-1\}\)，则（\quad）.
\choices{\(A\subseteq B\)}{\(\complement_{\mathbb R}A\subseteq\complement_{\mathbb R}B\)}{\(A\cap B=\varnothing\)}{\(A\cup B=\mathbb R\)}
\end{problem}

\begin{answer}
D
\end{answer}

\begin{solution}
由\(x^2-x-2=\left( x-2 \right)\left( x+1 \right)\ge 0\)，得到\(A=\left( -\infty,-1 \right]\cup\left[ 2,+\infty \right)\)，而\(B=\left( -1,+\infty \right)\). \(-1\in A\)但\(-1\notin B\)，所以\(A\subseteq B\)不成立，排除A. \(\complement_{\R}A=\left( -1,2 \right)\)，\(\complement_{\R}B=\left( -\infty,-1 \right]\)，例如\(0\)属于前者但不属于后者，所以第二项包含关系不成立，排除B. \(A\cap B=\left[ 2,+\infty \right)\ne\varnothing\)，排除C. \(A\cup B=\left( -\infty,-1 \right]\cup\left( -1,+\infty \right)=\R\)，故选D.
\end{solution}

\begin{problem}
已知函数 \(y=f(x)\) 的定义域为 \(\mathbb R\)，下列是 \(f(x)\) 无最大值的充分条件是（\quad）.
\choices
    {\(f(x)\) 为偶函数且关于直线 \(x=1\) 对称}
    {\(f(x)\) 为偶函数且关于点 \((1,1)\) 对称}
    {\(f(x)\) 为奇函数且关于直线 \(x=1\) 对称}
    {\(f(x)\) 为奇函数且关于点 \((1,1)\) 对称}
\end{problem}

\begin{answer}
D
\end{answer}

\begin{solution}
逐项判断如下. A项中，偶函数指\(f\left( -x \right)=f\left( x \right)\)，关于直线\(x=1\)对称指\(f\left( 1+t \right)=f\left( 1-t \right)\)，此时\(f\left( x \right)=\cos\left( \pi x \right)\)满足两个条件且在\(x=0\)处取得最大值\(1\)，因此A不是无最大值的充分条件，排除A. B项中，偶函数结合关于点\(\left( 1,1 \right)\)对称指\(f\left( 1+t \right)+f\left( 1-t \right)=2\)，此时常值函数\(f\left( x \right)=1\)满足两个条件且取得最大值\(1\)，因此B不是充分条件，排除B. C项中，奇函数结合关于直线\(x=1\)对称，此时零函数\(f\left( x \right)=0\)满足两个条件且取得最大值\(0\)，因此C不是充分条件，排除C. D项中，奇函数给出\(f\left( 0 \right)=0\)，点对称给出\(f\left( 1+t \right)+f\left( 1-t \right)=2\)对一切实数\(t\)成立. 在该恒等式中用奇函数改写第二项为\(-f\left( t-1 \right)\)，得到\(f\left( t+1 \right)-f\left( t-1 \right)=2\)，令\(u=t-1\)就有\(f\left( u+2 \right)=f\left( u \right)+2\)对一切实数\(u\)成立. 取\(u=0\)并归纳就有\(f\left( 2k \right)=2k\)对一切正整数\(k\)成立，函数值沿偶数点趋于正无穷，不存在上界，更不存在最大值. 于是D是无最大值的充分条件，故选D.
\end{solution}

\begin{problem}
在 \(\triangle ABC\) 中，\(D\) 为 \(BC\) 中点，\(E\) 为 \(AD\) 中点. 判断：
\begin{circlelist}
\item 存在 \(\triangle ABC\)，使得 \(\overrightarrow{AB}\cdot\overrightarrow{CE}=0\)；
\item 存在 \(\triangle ABC\)，使得 \(\overrightarrow{CE}\parallel(\overrightarrow{CB}+\overrightarrow{CA})\).
\end{circlelist}
则（\quad）.
\choices{\(\circled{1}\) 成立，\(\circled{2}\) 成立}{\(\circled{1}\) 成立，\(\circled{2}\) 不成立}{\(\circled{1}\) 不成立，\(\circled{2}\) 成立}{\(\circled{1}\) 不成立，\(\circled{2}\) 不成立}
\end{problem}

\begin{answer}
B
\end{answer}

\begin{solution}
以\(C\)为原点记位置向量为\(\bs{a}\)，\(\bs{b}\)的点为\(A\)，\(B\)，三角形非退化等价于\(\bs{a}\)与\(\bs{b}\)线性无关. 此时\(D=\dfrac{\bs{b}}{2}\)，\(E=\dfrac{\bs{a}}{2}+\dfrac{\bs{b}}{4}\)，\(\overrightarrow{CE}=\dfrac{\bs{a}}{2}+\dfrac{\bs{b}}{4}\)，\(\overrightarrow{CB}+\overrightarrow{CA}=\bs{a}+\bs{b}\). 若其中之一为零向量，则\(\bs{a}\)与\(\bs{b}\)线性相关，三角形退化，因此在真正的三角形中两者均为非零向量，平行等价于存在非零实数\(\mu\)使\(\dfrac{\bs{a}}{2}+\dfrac{\bs{b}}{4}=\mu\left( \bs{a}+\bs{b} \right)\). 由线性无关比较系数必须同时有\(\mu=\dfrac{1}{2}\)与\(\mu=\dfrac{1}{4}\)，不可能成立，所以\(\overrightarrow{CE}\)与\(\overrightarrow{CB}+\overrightarrow{CA}\)在任何三角形中恒不平行，\circled{2}不成立. 另一方面，取坐标\(A=\left( 0,0 \right)\)，\(B=\left( 2,0 \right)\)，\(C=\left( \dfrac{2}{3},1 \right)\)，则\(D=\left( \dfrac{4}{3},\dfrac{1}{2} \right)\)，\(E=\left( \dfrac{2}{3},\dfrac{1}{4} \right)\)，\(\overrightarrow{AB}=\left( 2,0 \right)\)，\(\overrightarrow{CE}=\left( 0,-\dfrac{3}{4} \right)\)，内积为\(0\)，三点不共线构成真正的三角形，所以存在三角形使\(\overrightarrow{AB}\cdot\overrightarrow{CE}=0\)，\circled{1}成立. 综上，\circled{1}成立，\circled{2}不成立，故选B.
\end{solution}
\section{解答题}

\begin{problem}
四棱锥 \(P-ABCD\)，底面为正方形 \(ABCD\)，边长为 \(4\)，\(E\) 为 \(AB\) 中点，\(PE\perp\) 平面 \(ABCD\).
\begin{enumerate}
\item 若 \(\triangle PAB\) 为等边三角形，求四棱锥 \(P-ABCD\) 的体积；
\item 若 \(CD\) 的中点为 \(F\)，\(PF\) 与平面 \(ABCD\) 所成角为 \(45^\circ\)，求 \(PD\) 与 \(AC\) 所成角的大小.
\end{enumerate}
\end{problem}

\begin{solution}
\begin{enumerate}\item 因为\(AB=4\)，\(E\)为\(AB\)中点，\(\triangle PAB\)为等边三角形，所以中线\(PE=\sqrt{4^{2}-2^{2}}=2\sqrt{3}\). 因为\(PE\perp\)平面\(ABCD\)，所以\(PE\)即为四棱锥的高，底面面积为\(4^{2}=16\). 故\(V_{P-ABCD}=\frac{1}{3}\times 16\times 2\sqrt{3}=\frac{32\sqrt{3}}{3}\). \item 因为\(PE\perp\)平面\(ABCD\)，所以\(EF\)为\(PF\)在底面的投影，\(PF\)与底面所成角即为\(\angle PFE\). 由\(\angle PFE=45^{\circ}\)得\(PE=EF\). 在正方形\(ABCD\)中，\(E\)，\(F\)分别为\(AB\)，\(CD\)中点，故\(EF=AD=4\)，从而\(PE=4\). 以\(E\)为原点，直线\(EB\)为\(x\)轴，直线\(EF\)为\(y\)轴，直线\(EP\)为\(z\)轴建立空间直角坐标系，则\(E(0,0,0)\)，\(A(-2,0,0)\)，\(B(2,0,0)\)，\(C(2,4,0)\)，\(D(-2,4,0)\)，\(P(0,0,4)\). 于是\(\overrightarrow{PD}=(-2,4,-4)\)，\(\overrightarrow{AC}=(4,4,0)\). \(|\overrightarrow{PD}|=\sqrt{4+16+16}=6\)，\(|\overrightarrow{AC}|=\sqrt{16+16}=4\sqrt{2}\)，\(\overrightarrow{PD}\cdot\overrightarrow{AC}=-8+16+0=8\). 记异面直线\(PD\)与\(AC\)所成角为\(\theta\)，则\(\cos\theta=\frac{|\overrightarrow{PD}\cdot\overrightarrow{AC}|}{|\overrightarrow{PD}||\overrightarrow{AC}|}=\frac{8}{6\times 4\sqrt{2}}=\frac{\sqrt{2}}{6}\). 故所成角的大小为\(\arccos\frac{\sqrt{2}}{6}\). \end{enumerate}
\end{solution}

\begin{problem}
已知 \(A\)，\(B\)，\(C\) 为 \(\triangle ABC\) 的三个内角，\(a\)，\(b\)，\(c\) 是其三条边，\(a=2\)，\(\cos C=-\dfrac14\).
\begin{enumerate}
\item 若 \(\sin A=2\sin B\)，求 \(b\)，\(c\)；
\item 若 \(\cos\left(A-\dfrac{\pi}{4}\right)=\dfrac45\)，求 \(c\).
\end{enumerate}
\end{problem}

\begin{solution}
\begin{enumerate}\item 由正弦定理，\(\sin A=2\sin B\)等价于\(a=2b\). 因为\(a=2\)，所以\(b=1\). 由余弦定理\(\cos C=\frac{a^{2}+b^{2}-c^{2}}{2ab}\)，代入\(a=2\)，\(b=1\)，\(\cos C=-\frac{1}{4}\)得\(\frac{4+1-c^{2}}{4}=-\frac{1}{4}\)，解得\(c^{2}=6\). 因为\(c>0\)，所以\(c=\sqrt{6}\). \item 因为\(\cos C=-\frac{1}{4}\)，\(0<C<\pi\)，所以\(\sin C=\sqrt{1-\cos^{2}C}=\frac{\sqrt{15}}{4}\)，且\(C\)为钝角，故\(c\)为最大边，\(c>a\)，由正弦定理\(\sin C>\sin A\). 由\(\cos\left(A-\frac{\pi}{4}\right)=\frac{4}{5}\)得\(\frac{\sqrt{2}}{2}(\cos A+\sin A)=\frac{4}{5}\)，即\(\cos A+\sin A=\frac{4\sqrt{2}}{5}\). 与\(\cos^{2}A+\sin^{2}A=1\)联立，记\(s=\sin A\)，则\(s^{2}+\left(\frac{4\sqrt{2}}{5}-s\right)^{2}=1\)，化简得\(2s^{2}-\frac{8\sqrt{2}}{5}s+\frac{7}{25}=0\)，解得\(s=\frac{\sqrt{2}}{10}\)或\(s=\frac{7\sqrt{2}}{10}\). 因为\(\left(\frac{7\sqrt{2}}{10}\right)^{2}=\frac{98}{100}>\frac{15}{16}=\left(\frac{\sqrt{15}}{4}\right)^{2}\)，即\(\frac{7\sqrt{2}}{10}>\sin C\)，与\(\sin A<\sin C\)矛盾，舍去. 故\(\sin A=\frac{\sqrt{2}}{10}\). 由正弦定理\(\frac{a}{\sin A}=\frac{c}{\sin C}\)得\(c=a\cdot\frac{\sin C}{\sin A}=2\times\frac{\sqrt{15}}{4}\times\frac{10}{\sqrt{2}}=\frac{5\sqrt{30}}{2}\). \end{enumerate}
\end{solution}

\begin{problem}
团队在 \(O\) 点西侧，东侧 \(20\) 千米处设有 \(A\)，\(B\) 两站点，测量一点 \(P\) 满足 \(|PA|-|PB|=20\) 千米. 以 \(O\) 为原点，东侧为 \(x\) 轴正半轴，北侧为 \(y\) 轴正半轴建立平面直角坐标系，\(P\) 在北偏东 \(60^\circ\) 处.
\begin{enumerate}
\item 求双曲线标准方程和 \(P\) 点坐标；
\item 又在南侧，北侧 \(15\) 千米处设 \(C\)，\(D\) 两站点，若 \(|QA|-|QB|=30\) 千米，\(|QC|-|QD|=10\) 千米，求 \(|OQ|\) 和 \(Q\) 点位置.
\end{enumerate}
\end{problem}

\begin{solution}
\begin{enumerate}\item 由\(|PA|-|PB|=20\)知\(P\)在以\(A(-20,0)\)，\(B(20,0)\)为焦点的双曲线右支上，故\(2a=20\)，\(c=20\)，得\(a=10\)，\(b^{2}=c^{2}-a^{2}=300\). 标准方程为\(\frac{x^{2}}{100}-\frac{y^{2}}{300}=1\)，\(x>0\). 北偏东\(60^{\circ}\)即从\(y\)轴正向向\(x\)轴正向偏\(60^{\circ}\)，与\(x\)轴夹角为\(30^{\circ}\)，射线\(OP\)为\(y=\tan 30^{\circ}x=\frac{\sqrt{3}}{3}x\)，\(x\ge 0\). 代入双曲线得\(\frac{x^{2}}{100}-\frac{x^{2}}{900}=1\)，即\(\frac{2x^{2}}{225}=1\)，得\(x=\frac{15\sqrt{2}}{2}\)，\(y=\frac{5\sqrt{6}}{2}\). 故\(P\left(\frac{15\sqrt{2}}{2},\frac{5\sqrt{6}}{2}\right)\). \item 由\(|QA|-|QB|=30\)知\(2a=30\)，\(c=20\)，得\(a=15\)，\(b^{2}=400-225=175\)，右支方程为\(\frac{x^{2}}{225}-\frac{y^{2}}{175}=1\)，\(x>0\). 设\(C(0,-15)\)，\(D(0,15)\)，由\(|QC|-|QD|=10>0\)知\(Q\)靠近\(D\)即北侧，\(2a=10\)，\(c=15\)，得\(a=5\)，\(b^{2}=225-25=200\)，上支方程为\(\frac{y^{2}}{25}-\frac{x^{2}}{200}=1\)，\(y>0\). 联立得\(x^{2}=225+\frac{9y^{2}}{7}\)，代入第二式得\(y^{2}\left(\frac{1}{25}-\frac{9}{1400}\right)=1+\frac{225}{200}\)，即\(y^{2}\cdot\frac{47}{1400}=\frac{17}{8}\)，故\(y^{2}=\frac{2975}{47}\)，\(x^{2}=\frac{14400}{47}\). 取正值得\(Q\left(\sqrt{\frac{14400}{47}},\sqrt{\frac{2975}{47}}\right)\). 于是\(|OQ|=\sqrt{x^{2}+y^{2}}=\sqrt{\frac{17375}{47}}\)千米\(\approx 19.23\)千米\(\approx 19227\)米. 记\(OQ\)与\(x\)轴夹角为\(\theta\)，则\(\tan\theta=\frac{y}{x}=\sqrt{\frac{2975}{14400}}\approx 0.4545\)，\(\theta\approx 24^{\circ}\)，故北偏东\(90^{\circ}-\theta\approx 66^{\circ}\). 按原PDF精确到1米、1°计，\(|OQ|\approx 19227\)米，\(Q\)在北偏东约\(66^{\circ}\)方向. 原PDF第4页要求\(|OQ|\)精确到1米、\(Q\)位置精确到1米与\(1^{\circ}\)，本页现稿略去精度要求但距离条件无出入，以上精确值与近似值均为原PDF题面所得. \end{enumerate}
\end{solution}

\begin{problem}
已知函数 \(f(x)=|x+a|-a-x\).
\begin{enumerate}
\item 若 \(a=1\)，求函数定义域；
\item 若 \(a\ne0\)，且 \(f(ax)=a\) 有 \(2\) 个不同实数根，求 \(a\) 的取值范围；
\item 是否存在实数 \(a\)，使得函数 \(f(x)\) 在定义域内具有单调性？若存在，求出 \(a\) 的取值范围.
\end{enumerate}
\end{problem}

\begin{solution}
\begin{enumerate}\item 按原PDF第4页题面\(f(x)=\sqrt{|x+a|-a}-x\)求解如下（本页现稿漏写根号，写作\(|x+a|-a-x\)；若按现稿字面函数则定义域为\(\R\)，结论不同，以下均为原PDF题面所得）.  \(f(x)=\sqrt{|x+a|-a}-x\). 当\(a=1\)时被开方数\(|x+1|-1\ge 0\)，即\(|x+1|\ge 1\)，解得\(x\le -2\)或\(x\ge 0\). 故定义域为\((-\infty,-2]\cup[0,+\infty)\). \item 设\(a\ne 0\)，\(f(ax)=a\)即\(\sqrt{|ax+a|-a}-ax=a\)，等价于\(\sqrt{|ax+a|-a}=ax+a\)且\(ax+a\ge 0\). 记\(t=ax+a\ge 0\)，则\(|t|=t\)，条件化为\(t-a=t^{2}\)，即\(a=t-t^{2}\)，\(t\ge 0\). 因为\(a\ne 0\)，\(x=\frac{t}{a}-1\)与\(t\)一一对应，故\(x\)有两个不同实根当且仅当\(a=t-t^{2}\)，\(t\ge 0\)有两个不同实根. 记\(g(t)=t-t^{2}=\frac{1}{4}-\left(t-\frac{1}{2}\right)^{2}\)，\(t\ge 0\). 若\(a>\frac{1}{4}\)则无解；若\(a=\frac{1}{4}\)则\(t=\frac{1}{2}\)唯一；若\(a=0\)舍去（且题设\(a\ne 0\)）；若\(a<0\)则\(t^{2}-t+a=0\)判别式大于零，两根之积为\(a<0\)，一正一负，在\(t\ge 0\)中恰一根. 若\(0<a<\frac{1}{4}\)则判别式\(1-4a>0\)，两根为\(\frac{1\pm\sqrt{1-4a}}{2}\)，之和为\(1\)，之积为\(a>0\)，故两根均为正数且互异. 此时对应两个不同的\(x\)，且均满足\(t-a=t^{2}\ge 0\)与\(t\ge 0\)，回代成立. 故\(a\)的取值范围为\(\left(0,\frac{1}{4}\right)\). \item 先定定义域：需\(|x+a|\ge a\). 若\(a\le 0\)则恒成立，定义域为\(\R\)；若\(a>0\)则\(x\ge 0\)或\(x\le -2a\)，定义域为\((-\infty,-2a]\cup[0,+\infty)\). 当\(x\ge -a\)时\(|x+a|=x+a\)，\(f(x)=\sqrt{x}-x\)（需\(x\ge 0\)）；当\(x<-a\)时\(f(x)=\sqrt{-x-2a}-x\). 在\(x<-a\)一侧，\(-x-2a\)随\(x\)增大而减小，故\(\sqrt{-x-2a}-x\)严格递减（有定义处）. 在\(x\ge -a\)一侧，\(f(x)=\sqrt{x}-x\)，对\(x_{1},x_{2}\ge \frac{1}{4}\)，\(\sqrt{x_{2}}-\sqrt{x_{1}}=\frac{x_{2}-x_{1}}{\sqrt{x_{2}}+\sqrt{x_{1}}}<x_{2}-x_{1}\)，故在\([\frac{1}{4},+\infty)\)上严格递减，而在\([0,\frac{1}{4}]\)上严格递增. 若\(a\le -\frac{1}{4}\)，则\(-a\ge \frac{1}{4}\)，右侧区间\([-a,+\infty)\)落在递减段，故\(f\)在\([-a,+\infty)\)与\((-\infty,-a)\)上均严格递减；又\(f\)在\(x=-a\)处左右均为\(\sqrt{-a}+a\)，连续，故\(f\)在\(\R\)上严格递减，具有单调性. 若\(a>0\)，则\(0,\frac{1}{4},1\)均在定义域右支上，\(f(0)=0\)，\(f\left(\frac{1}{4}\right)=\frac{1}{4}\)，\(f(1)=0\)，既有递增对又有递减对，不单调. 若\(-\frac{1}{4}<a\le 0\)，则\(-a\in[0,\frac{1}{4})\)，右侧含\(-a,\frac{1}{4},1\)，\(f(-a)=\sqrt{-a}+a=\frac{1}{4}-\left(\sqrt{-a}-\frac{1}{2}\right)^{2}<\frac{1}{4}=f\left(\frac{1}{4}\right)\)，而\(f\left(\frac{1}{4}\right)>f(1)\)，故不单调. 因此存在，且\(a\)的取值范围为\(\left(-\infty,-\frac{1}{4}\right]\). 以上三问结论均为原PDF含根号题面所得；按本页现稿无根号字面则结论不同，不采用. \end{enumerate}
\end{solution}

\begin{problem}
已知数列 \(\{a_n\}\) 满足 \(a_n\ge0\)，对任意 \(n\ge2\)，\(a_n\) 和 \(a_{n+1}\) 中存在一项使其为另一项与 \(a_{n-1}\) 的等差中项.
\begin{enumerate}
\item 已知 \(a_1=5\)，\(a_2=3\)，\(a_4=2\)，求 \(a_3\) 的所有可能取值；
\item 已知 \(a_1=a_4=a_7=0\)，\(a_2\)，\(a_5\)，\(a_8\) 为正数，证明：\(a_2\)，\(a_5\)，\(a_8\) 成等比数列，并求公比 \(q\)；
\item 已知数列中恰有 \(3\) 项为 \(0\)，即 \(a_r=a_s=a_t=0\)，\(2<r<s<t\)，且 \(a_1=1\)，\(a_2=2\)，求 \(a_{r+1}+a_{s+1}+a_{t+1}\) 的最大值.
\end{enumerate}
\end{problem}

\begin{solution}
\begin{enumerate}\item 对\(n\ge 2\)，条件等价于\(2a_{n}=a_{n+1}+a_{n-1}\)或\(2a_{n+1}=a_{n}+a_{n-1}\)，即差分\(d_{n}=a_{n+1}-a_{n}\)满足\(d_{n}=d_{n-1}\)或\(d_{n}=-\frac{d_{n-1}}{2}\). 当\(n=2\)时，\(2a_{2}=a_{3}+a_{1}\)或\(2a_{3}=a_{2}+a_{1}\)，代入\(a_{1}=5\)，\(a_{2}=3\)得\(a_{3}=1\)或\(a_{3}=4\). 当\(n=3\)时，\(2a_{3}=a_{4}+a_{2}\)或\(2a_{4}=a_{3}+a_{2}\)，代入\(a_{2}=3\)，\(a_{4}=2\)得\(a_{3}=\frac{5}{2}\)或\(a_{3}=1\). \(a_{3}\)须同时满足两处，交集为\(a_{3}=1\). 回代检验：\(2a_{2}=6=a_{3}+a_{1}\)，\(2a_{4}=4=a_{3}+a_{2}\)，成立且非负. 故\(a_{3}\)的唯一可能取值为\(1\). \item 设\(b=a_{2}>0\). \(n=2\)时\(a_{3}=2b\)或\(a_{3}=\frac{b}{2}\). 若\(a_{3}=2b\)，则\(n=3\)处\(a_{2}=b\)，\(a_{3}=2b\)，\(a_{4}=0\)须满足\(4b=b\)或\(0=3b\)，均与\(b>0\)矛盾，故\(a_{3}=\frac{b}{2}\). 此时\(n=3\)处\(2a_{3}=b=a_{4}+a_{2}\)成立，另一分支\(0=\frac{3b}{2}\)不成立. \(n=4\)处涉\(a_{3}=\frac{b}{2}\)，\(a_{4}=0\)，\(a_{5}>0\)：分支\(2a_{4}=a_{5}+a_{3}\)给出\(a_{5}=-\frac{b}{2}<0\)，舍去，故必\(2a_{5}=a_{4}+a_{3}\)，得\(a_{5}=\frac{b}{4}\). \(n=5\)处\(a_{5}=\frac{b}{4}\)，分支为\(a_{6}=\frac{b}{2}\)或\(a_{6}=\frac{b}{8}\). 若\(a_{6}=\frac{b}{2}\)，则\(n=6\)处\(a_{5}=\frac{b}{4}\)，\(a_{6}=\frac{b}{2}\)，\(a_{7}=0\)须满足\(b=\frac{b}{4}\)或\(0=\frac{3b}{4}\)，均矛盾，故\(a_{6}=\frac{b}{8}\)，且\(n=5\)取第二分支成立. \(n=6\)处\(2a_{6}=\frac{b}{4}=a_{7}+a_{5}\)成立，另一分支不成立. \(n=7\)处涉\(a_{6}=\frac{b}{8}\)，\(a_{7}=0\)，\(a_{8}>0\)：分支\(2a_{7}=a_{8}+a_{6}\)给出负值，舍去，故\(2a_{8}=a_{7}+a_{6}\)，得\(a_{8}=\frac{b}{16}\). 于是\(a_{2}=b\)，\(a_{5}=\frac{b}{4}\)，\(a_{8}=\frac{b}{16}\)成等比数列，公比\(q=\frac{1}{4}\). \item 记\(d_{n}=a_{n+1}-a_{n}\)，则\(d_{n}=d_{n-1}\)或\(d_{n}=-d_{n-1}/2\)，\(d_{1}=1\). 若某处出现连续两项为零，则向前向后递推均迫使全列为零，与\(a_{1}=1\)矛盾，故在恰3项为零时零点皆孤立. 设\(a_{r}=0\)，\(r>2\)，则\(n=r\)处分支\(2a_{r}=a_{r+1}+a_{r-1}\)将给出\(a_{r+1}+a_{r-1}=0\)，即\(a_{r-1}=0\)，形成连续零点，不可能，故必取\(2a_{r+1}=a_{r}+a_{r-1}\)，得\(a_{r+1}=\frac{a_{r-1}}{2}>0\). 对\(n=r-1\)处作同样分支检验，分支\(2a_{r}=a_{r-1}+a_{r-2}\)将给出\(a_{r-1}=a_{r-2}=0\)，不可能，故必取\(2a_{r-1}=a_{r}+a_{r-2}\)，得\(a_{r-2}=2a_{r-1}\). 于是\(d_{r-1}=-a_{r-1}<0\)，\(d_{r}=a_{r+1}=\frac{a_{r-1}}{2}=-d_{r-1}/2\). 从\(d_{1}=1\)到\(d_{r-1}\)共\(r-2\)步，每步保持或乘\(-1/2\)，故\(d_{r-1}=\left(-\frac{1}{2}\right)^{i}\)，其中\(i\)为其中取第二分支的次数. 因为\(d_{r-1}<0\)，故\(i\)为奇数，\(a_{r-1}=\left(\frac{1}{2}\right)^{i}\)，\(a_{r+1}=\left(\frac{1}{2}\right)^{i+1}\). 由\(i\le r-2\)且\(i\ge 1\)，最大值在\(i=1\)取得，为\(\frac{1}{4}\)，可达到（例如\(d_{2}=-\frac{1}{2}\)后保持至\(a_{6}=0\)）. 对第二个零点\(s\)，从\(d_{r}=\left(\frac{1}{2}\right)^{i+1}>0\)到\(d_{s-1}=-a_{s-1}<0\)需奇数次\(j\)取第二分支，得\(a_{s+1}=\left(\frac{1}{2}\right)^{i+j+2}\)，在\(i=j=1\)时最大为\(\frac{1}{16}\)，可达到. 对第三个零点\(t\)作同样分支检验，从\(d_{s}=\left(\frac{1}{2}\right)^{i+j+2}>0\)到\(d_{t-1}=-a_{t-1}<0\)需奇数次\(k\)取第二分支，得\(a_{t+1}=\left(\frac{1}{2}\right)^{i+j+k+3}\)，在\(i=j=k=1\)时最大为\(\frac{1}{64}\)，可达到（之后保持\(d\)为正不再归零，恰3个零点）. 故\(a_{r+1}+a_{s+1}+a_{t+1}\le\frac{1}{4}+\frac{1}{16}+\frac{1}{64}=\frac{21}{64}\).等号可达到，最大值为\(\frac{21}{64}\). \end{enumerate}
\end{solution}
